\section{Explicit output and one uniform program}\label{sec:uniformization}

We first reduce exact cardinality to subcubic perfect matching, retaining
all witness information. We then combine the fixed-parameter procedures
into the single program of Theorem~\ref{thm:main}. Throughout this section
$p=n+m$ refers to the original explicit input size.

\subsection{Witness-preserving reductions}

Read adjacency-list boundaries and neighbor occurrences once. Emit an edge
record, with its endpoints and a fresh identifier, only at the occurrence
$(u,v)$ with $u<v$. Input validity and simplicity imply that this produces
each edge exactly once, regardless of list order. Retain all vertex records,
including isolates. This takes $O(n+m)$ word operations, and the resulting
read-only data can be shared among branches. If $n=1$, return the empty list.

Two reductions connect the required output to the subcubic procedure.
First, a permutation network turns a cardinality threshold into a
perfect-matching instance. Second, a path gadget reduces the degree to
three. We give both constructions and their explicit witness decoders.
For a target cardinality $k$, adding $n-2k$ new vertices adjacent to every
original vertex would turn a size-$k$ matching into a perfect matching,
but could add quadratically many edges. The network replaces these
complete connections by $O(n\log(n+2))$ positions and arcs while allowing
any $n-2k$ original vertices to reach the required new vertices.

\begin{lemma}[A linearithmic permutation network]\label{uniform:network}
For every power of two $s\ge1$, there is a directed acyclic network with
$s$ distinct input positions and $s$ distinct output positions, of size
$O(s\log(s+2))$, in which every permutation of inputs to outputs is routable
by vertex-disjoint directed paths. A position may be both an input and an
output when $s=1$. The network can be constructed uniformly within its size
bound.
\end{lemma}

\begin{proof}
Use the recursive switching construction underlying permutation
networks~\cite{PermutationNetworks}. For $s=1$, take a single position and
its length-zero route. For $s>1$, take $s$ new input positions, $s$ new output
positions, and disjoint upper and lower networks on $s/2$ positions. Connect
each input in pair $2i-1,2i$ to input $i$ of both copies. Connect output $j$
of either copy to both positions in output pair $2j-1,2j$. The size recurrence
is $Q(s)=2Q(s/2)+O(s)$, and all arcs respect the resulting input-to-output
order.

Given a permutation, form a bipartite multigraph between its input pairs and
output pairs, with one edge per assigned input--output pair. Every vertex
has degree two, so the components are even cycles, allowing a cycle of two
parallel edges. Color each cycle alternately upper and lower. Every input
pair and output pair has one edge of each color. The upper edges specify a
permutation through the upper network, and the lower edges one through the
lower network. Route inductively and attach the corresponding outer arcs.
The two copies are disjoint and each outer position occurs once, giving
vertex-disjoint routes. The construction of the network itself uses only
the stated recursive lists and port identifiers.
\end{proof}

\begin{lemma}[Exact cardinality to perfect matching]\label{uniform:cardinality}
Given $G$ and an integer $0\le k\le\lfloor n/2\rfloor$, one can construct
in $O(m+n\log(n+2))$ time a loopless graph $H_k$ of that size such that
$H_k$ has a perfect matching if and only if $G$ has a matching of size $k$.
A perfect matching of $H_k$ yields an explicit size-$k$ matching of $G$
in linear time in the constructed size.
\end{lemma}

\begin{proof}
Put $b=n-2k$ and choose the least power of two $s\ge n$. Use the network of
Lemma~\ref{uniform:network} and select $b$ distinct output positions.
Retain all original vertices and edges. Split every network position $v$
into vertices $v_{\mathrm{in}},v_{\mathrm{out}}$ joined by an idle edge.
An arc $u\to v$ becomes the edge $u_{\mathrm{out}}v_{\mathrm{in}}$.
Join original vertex $i$ to the input endpoint of network input $i$.
Finally, add one mandatory sink vertex adjacent to the output endpoint
of each selected output. Every constructed vertex is required to be
covered by a perfect matching.

In any perfect matching, a network pair either uses its idle edge or has
both endpoints matched externally. The numbers of externally matched input
and output endpoints are therefore equal. Selected transition edges
contribute once to each number and cancel. Original-to-network edges
contribute one input endpoint, and edges to mandatory sinks contribute one
output endpoint. All $b$ sinks must be covered, so exactly $b$ original
vertices use network edges. The remaining $n-b=2k$ original vertices are
covered by original edges, which give the decoded matching.

Conversely, let a size-$k$ matching leave $b$ original vertices uncovered.
Map them bijectively to the selected outputs and extend to a permutation
of all $s$ positions. Route the permutation and keep just those $b$ paths.
On each retained path, select its original-to-input edge, all transition
edges, and the edge to its sink. Each active split position is covered once
at each endpoint. Use idle edges on all other positions and the original
matching on its covered vertices. This is a perfect matching of $H_k$.
For a length-zero route, the two outer edges cover the one split position.
For $b=0$, every network pair is idle.

The number of added positions and edges is $O(n\log(n+2))$. Store each edge's
type and original identifier. Decoding simply selects the original edges
from the returned perfect matching. No list ordering is assumed.
\end{proof}

The next reduction uses the path replacement for matching described by
Dahlhaus and Karpinski~\cite{DahlhausKarpinski}. We record the construction
and the witness conversion needed here.

\begin{lemma}[Degree reduction]\label{uniform:degree}
A loopless graph with $v$ vertices and $e$ edges either has an isolated
vertex, certifying absence of a perfect matching, or can be replaced in
$O(v+e)$ time by an equivalent perfect-matching instance of maximum degree
three and size $O(v+e)$. Witnesses can be converted in either direction
within the same bound.
\end{lemma}

\begin{proof}
Replace a vertex of degree $r\ge1$ by a path on $2r-1$ vertices, attaching
its $r$ old incidences to successive odd positions. Retain the original
edge identifiers at both ends. There are $r$ odd and $r-1$ even positions.
Every even position must be matched internally to an odd one. Exactly one
odd position therefore remains for an external edge. Hence each gadget
selects exactly one old incidence, and any perfect matching projects to
one in the original graph.

Conversely, the odd attachment of a selected original matching edge leaves
two even-length paths on its sides. Match those internally. This also
covers $r=1$, whose gadget is a single vertex. Each odd position has at
most two path neighbors and one external edge, and each even position
has degree two. The number of gadget vertices is
\[
                    \sum_u(2\deg(u)-1)=4e-v.
\]
A degree scan, incidence-port assignment, and path construction give the
time bound and the explicit conversions. Check for isolates first.
\end{proof}

Applying these reductions at a threshold $k>0$ gives a loopless subcubic
instance with
\begin{equation}\label{uniform:reduced-size}
                 N=O(m+n\log(n+2)),\qquad N\ge n,
\end{equation}
unless absence of a perfect matching is already certified by an isolate.
The latter inequality holds because every original vertex has at least
one gadget vertex. The construction on a simple input can itself use
simple graphs: network arcs, attachments and path edges have their
specified distinct positions. It is also covered by the intermediate
algorithms' individual-edge-copy convention.

For the perfect-matching procedure set $S$ to the first power of two at
least $16+n+m+N$. Thus
\begin{equation}\label{uniform:master}
               n\le S=O(p\log(p+2)).
\end{equation}
Check every proposed perfect matching by edge identifier and one-time
coverage of every gadget vertex. The reductions then decode its actual
original matching in the same auxiliary-size order.

\subsection{Threshold search and its error budget}

Each branch runs its own integer binary search. Initially the feasible
lower endpoint is zero, with the stored empty matching, and the exclusive
upper endpoint is $\lfloor n/2\rfloor+1$. An accepted test stores its
decoded witness and raises the lower endpoint; a negative test lowers the
upper endpoint. There are $O(\log(n+2))$, and certainly at most $n+1$,
tests. Exact-size feasibility is monotone because edges may be dropped
from a matching. Consequently, if no feasible test receives a false
negative, the final stored witness is maximum.

For each branch index $t\ge1$, use
\begin{equation}\label{uniform:branch-parameters}
                     a=1000+t,\qquad \sigma=t+20
\end{equation}
in Theorem~\ref{thm:subcubic}. A deterministic prerequisite suspension
suspends the entire branch with no output. A capped stochastic failure
terminates that perfect-matching test with an ordinary negative answer.
All invocations use fresh independent bits.

On a feasible perfect-matching instance, conditional on its input and
the earlier history, the probability of a false negative is at most
$2^{-\sigma}S^{-5}$. This bound does not condition on future
noninterruption. Each valid reached randomized invocation has its declared
conditional bound, while a deterministic suspension inside an invocation
does not itself issue a negative answer. The main guards permit the
coarse polynomial union bound in Theorem~\ref{thm:subcubic}; each backend
instance already includes its capped internal invocations in its
reliability estimate. There are no false positive perfect-matching answers.
By $S\ge n+1$ at all nontrivial tests and the bound on the number of
adaptive thresholds,
\begin{equation}\label{uniform:branch-error}
 \Pr[\text{branch $t$ completes with an incorrect matching}]
                         \le 2^{-(t+20)}.
\end{equation}
No guarantee of avoiding suspension on small inputs is needed for this
inequality.

Branch zero performs the same reduction and search, using
the deterministic procedure of Lemma~\ref{lem:fallback}. Every threshold
answer is exact, so this branch always completes with a maximum matching.
It uses a fixed polynomial number of instructions after finite-word
simulation. Indeed every simple input satisfies
$m\le n(n-1)/2$, and Equation~\eqref{uniform:reduced-size} makes every
fallback instance polynomial in $n+2$. No perfect-matching promise is used.

\subsection{Fixed branches and eventual completion}

\begin{lemma}[Every-path performance of a fixed branch]\label{uniform:fixed-branch}
For every sufficiently large fixed $t$, branch $t$ avoids deterministic
suspension on all sufficiently large inputs, on every computation path.
Its total instruction count, including unsuccessful tests, threshold
search, and explicit witness handling, is
\[
                          p^{1+O(1/a)+o_a(1)}.
\]
The coefficient in $O(1/a)$ is absolute.
\end{lemma}

\begin{proof}
Each test satisfies Equation~\eqref{uniform:master}. The guarded-prefix
estimates in Theorem~\ref{thm:subcubic} and Lemma~\ref{lem:flat-cycle}
show eventual adequacy at fixed sufficiently large $a,\sigma$, independently
of matchability and of which verified constructions or sampling failures
occurred. Every accepted sampled structure satisfies its checked
deterministic invariants, and rejected attempts obey their own unconditional caps.
The estimates therefore apply through a purported first deterministic
interruption. They do not require a favorable random stream. An exhausted
sampling routine terminates the test rather than suspending it.

The fixed-parameter time bound, the logarithmic number of thresholds,
and the reductions give the displayed exponent. Logarithmic and
fixed-parameter subpolynomial factors absorb the substitution
$S=O(p\log(p+2))$.

All parameters and adequacy checks are computed by finite integer loops:
bit lengths, rounded powers of two, products, quotients, and comparisons
of scaled dyadics or ratios. Inadequate giant parameters need not pass
their checks. The assertion that fixed-parameter logarithmic factors
eventually fit size, multiplicity or quality margins requires no encoded
table of thresholds. Fixed-parameter bit lengths and arithmetic/dictionary
simulation costs are logarithmic powers, or are already charged in the
static flow primitive. Individual random-bit costs are included.
\end{proof}

\subsection{An instruction-level schedule}

Weighted dovetailing, familiar from universal
search~\cite{LevinUniversal}, lets one program compete with every fixed
parameter branch. The proof also accounts for summable error allowances
and finite-word execution.

Run rounds $j=0,1,2,\ldots$. In round $j$, give branches
$t=0,\ldots,j$, in that order, at most $2^{j-t}$ elementary implementation
steps each, resuming a branch at its interruption point. A suspended branch
stays inactive. Stop when the first branch completes with its answer.
The winning branch constructs its decoded output privately within its
charged work; copying the list of at most $n/2$ original labeled edges
adds input-linear work.

Here a scheduled step means an actual bounded word-RAM instruction with
constant control overhead, not a whole graph routine, arbitrary-precision
operation, or allocation request. The parameter algorithms consist of a
fixed program of graph, list, tree and integer subroutines, the one fixed
fallback, and the one fixed integral-flow implementation. Store each
branch as a coroutine with a program counter and stack. Use charged
digit arrays, even unit-bit digits if desired, for longer logical values.
Long virtual indices and addresses are keys in dictionaries of touched
cells with a default initial value. Comparisons, initialization, recursion,
and dictionary access are themselves interruptible implementations.

Allocate and initialize actual records one word at a time. A paused
computation retains its state without copying it; its physical storage
is proportional to actual executed allocations plus the shared input.
Arenas may retain already spent cells. In particular, constructing an
enormous power for a large varying parameter is charged and can pause
inside its arithmetic or allocation loop. It is not an atomic reservation
of an enormous memory interval. For each fixed branch of
Lemma~\ref{uniform:fixed-branch}, these implementations add only the
logarithmic-power overheads already allowed. The fixed fallback simulation
still has a polynomial bound.

\begin{lemma}[Halting and one physical word constant]\label{uniform:physical-ram}
The schedule is realizable by one finite program with a fixed physical word
constant $c_0$. On every path, including failed paths, it halts after
polynomially many instructions and uses polynomially many addressed words,
with constants independent of the input.
\end{lemma}

\begin{proof}
First ignore physical-address overflow, and fix the fallback's finite
digit-array and touched-cell simulation. Its logical clock, bit lengths,
and dictionary costs give a fixed polynomial bound $F(n)$ on its actual
work and records. This estimate can use unit-bit data and physical pointer
operations; it makes no assumption that the varying parameters of other
branches fit a word.

Unless the schedule has already stopped, the fallback receives $2^j$ steps
in round $j$. It therefore finishes within $J=O(\log(n+2))$ rounds, with
a fixed constant. Total branch allowances through round $J$ are
\[
 \sum_{j=0}^J\sum_{t=0}^j2^{j-t}<2^{J+2}.
\]
Starting and switching branches, status checks, and loop control add only
polynomial logarithmic overhead; input/setup and output are polynomial
as well. Since every allocation is charged, all actually touched physical
records before halt are bounded by a fixed polynomial in $n+2$, even
if every arena retains spent cells.

Use independent logical indexing in each branch and actual physical
addresses only as pointers, for example with sequential allocation.
Interleaving cannot change a simulated algorithm. Choose $c_0$ once,
large enough that this polynomial bound covers physical addresses,
quota counters, program control, and input indices in the prescribed
$w$-bit words. No overflow then occurs on any path. Long logical integers
remain digit arrays. Smaller logical words can be masked or simulated;
long shifts, divisions and other operations use the corresponding charged
loops. Thus the choice of $c_0$ follows the fixed polynomial bound rather
than being a premise about unbounded logical arithmetic.
\end{proof}

For the probability calculation, couple the schedule to independent
infinite fair-bit streams for its branches, revealing bits only when
requested. This is equivalent to generating fresh charged bits during
execution. Consider each branch's own complete computation on its stream,
including a possible suspension. Interleaving and the identity of the
first finisher impose no conditioning on the error estimates.
Branch zero is exact, so Equation~\eqref{uniform:branch-error} gives
\begin{equation}\label{uniform:total-error}
 \Pr[\text{the scheduled output is incorrect}]
 \le\sum_{t\ge1}2^{-(t+20)}<\frac13.
\end{equation}

\subsection{A common exponent envelope}

\begin{proof}[Proof of Theorem~\ref{thm:main}]
Halting and the computation model follow from
Lemma~\ref{uniform:physical-ram}; explicit output and success follow
from the reductions, threshold search, and
Equation~\eqref{uniform:total-error}.

Fix a sufficiently large branch index $t$. If it needs $T$ steps,
round $j=t+\lceil\log_2\max(1,T)\rceil$ alone gives it enough quota,
unless another answer has already terminated the schedule. The total
budget through that round is at most a constant times $2^t(T+1)$,
in addition to input and logarithmic control overhead. Thus the single
schedule competes with every fixed branch with a fixed-index factor.
For every $\xi>0$, choose $t$ large enough in the analysis that the
absolute $O(1/a)$ exponent in Lemma~\ref{uniform:fixed-branch} is smaller
than a suitable fraction of $\xi$. For all sufficiently large $p$, its
$o_a(1)$ term, constant factor and logarithmic overhead fit the remaining
slack. Consequently the \emph{same program} has every-path time at most
$p^{1+\xi}$ for all sufficiently large inputs. No choice of $t$ based on
the input is made by the program.

To exhibit one envelope, let $T(G,\omega)$ be the actual total instruction
count of the schedule on input $G$ and random path $\omega$. For $x\ge2$
define
\[
 e(x)=\sup_{\substack{G,\omega:\,n(G)+m(G)\ge x}}
 \max\left\{0,
 \log_{n(G)+m(G)}\bigl(\max\{1,T(G,\omega)\}\bigr)-1\right\}.
\]
The polynomial fallback bounds this supremum, and the preceding
every-path estimates imply $e(x)\to0$. Put
$\eta(q)=e(\max\{2,q\})$ for integers $q\ge1$. This function is
nonnegative and nonincreasing, and it tends to zero. If $p\ge2$, the term
for $(G,\omega)$ in the supremum defining $e(p)$ gives
\[
                    T(G,\omega)\le p^{1+\eta(p)}.
\]
The only input with $p=1$ is the one-vertex empty graph, which is handled
directly. Choose a fixed $C\ge1$ large enough to cover its instruction count.
The all-input bound in Theorem~\ref{thm:main} follows.

For the sparse specialization, let $E=\sup_{q\ge1}\eta(q)<\infty$.
If $m\le\kappa n$, then $n\le p\le(1+\kappa)n$ and monotonicity gives
\[
 T(G,\omega)\le C p^{1+\eta(p)}
 \le C(1+\kappa)^{1+E}n^{1+\eta(n)}.
\]
Thus the same algorithm and envelope apply for every fixed $\kappa$.
When $m=\Theta(n^2)$, one instead has $p=\Theta(n^2)$, giving the stated
$n^{2+o(1)}$ dense specialization.

Every operation, random bit, intermediate construction and output
copy is part of $T$, so no further uncharged term remains. The input
conversion did not assume connectedness, nonisolation, or a particular
list order. This proves all assertions of Theorem~\ref{thm:main}.
\end{proof}
